\subsection{A common-height estimate before a boundary root}
\label{sec:direct-quadratic-prefix}

Beyond the contraction range, suppose both drifts still have positive
margins. Equal leading actions then relate the two radii at the same
central height. The derivative of this relation has a useful
cancellation: the quadratic coefficients drop out of the derivative
ratio. This determines the crossing direction of \(\mathcal L\)
before the positive branch reaches a boundary root, for either sign
of the quadratic coefficient.

Suppose the scales \eqref{eq:direct-quadratic-prefix-scales} hold with
\(\delta\) and \(\epsilon_3=\ell/\delta\) sufficiently small. The drifts
\(K_\sigma\) of \cref{sec:early-central-heights} have the quadratic
comparisons
\[
 K_\sigma^0(x)=a-\sigma bx+\bar c x^2.
\]
Choose a positive terminal radius \(R\), possibly varying
over the compact coefficient set, such that \(K_+^0\)
has a fixed positive margin on \(0\le x\le R\).
Assume also that the equal-action negative radii for these
terminal actions lie in a fixed compact interval on which
\(K_-^0\) has a fixed positive margin.

On these comparison intervals set
\[
 \mathcal A_\sigma^0(x)=\int_0^x\frac{u\,du}{K_\sigma^0(u)},\qquad
 \mathcal A_-^0(y)=\mathcal A_+^0(x),\qquad y=f(x).
\]
The drift margins make \(\mathcal A_-^0\) strictly increasing, so its
inverse at the prescribed actions defines \(f\) on the chosen compact
intervals.

\begin{samepage}
\begin{lemma}[Equal-action comparison]
\label{lem:direct-equal-action-comparison}
For \(x,y=f(x)>0\) in these intervals, put
\(R_1(x)=K_-^0(y)/K_+^0(x)\). Then
\begin{equation}
\label{eq:direct-quadratic-critical-identity}
 \frac{yy'}x=R_1,\qquad
 \frac{R_1'}{R_1}
       =\frac{b(1+x/y)}{K_+^0(x)},\qquad
 \frac{(D+x^2-\lambda y^2)'}{2x}=1-\lambda R_1
\end{equation}
for any constants \(D,\lambda\). In particular, \(R_1\) is
nondecreasing when \(b\ge0\).
\end{lemma}

\begin{proof}
Differentiate \(\mathcal A_-^0(y)=\mathcal A_+^0(x)\) to obtain the first identity.
Since \((K_-^0)'(y)=b+2\bar c y\) and
\((K_+^0)'(x)=-b+2\bar c x\), logarithmic differentiation
of \(R_1\) and substitution of the first identity cancel
the two quadratic terms and give the second. The third
follows by differentiating \(D+x^2-\lambda y^2\).
\end{proof}
\end{samepage}

\begin{lemma}[Quadratic prefix]
\label{lem:direct-quadratic-prefix}
Fix a small \(d_f>0\) and put \(z_e=\delta\ell^2d_f\).
On the central heights in \(\cH_p\) satisfying
\[
 h\ge z_e,\qquad\hbox{or}\qquad
 0<h<z_e,\quad |\eta_+(z_e;h)|/\ell\le R,
\]
the return displacement has at most two distinct zeros.
Moreover, for any logarithmic height coordinate \(T=-c_T\log h\),
where \(c_T>0\)
is independent of the central height, every zero of
\[
 \mathcal L(T)=\log\!\left[
       \Pi_p'(I_-(h))\frac{I_{-,h}(h)}{I_{+,h}(h)}\right]
\]
on this interval satisfies \(\mathcal L_T>0\).
The constants are uniform on the compact coefficient and
action domains just specified.
\end{lemma}

\begin{proof}
By \cref{lem:early-central-heights}, there is \(r_0\in(0,R)\) such that
the central heights with \(h\ge z_e\), or with \(h<z_e\) and
\(x_+=|\eta_+(z_e;h)|/\ell\le r_0\), form an initial interval in \(T\) on
which the return contracts; in particular \(\mathcal L<0\) there.

\emph{Matching the branches.}
It remains to consider \(r_0\le x=x_+\le R\). On the prescribed
compact drift intervals, the central-height identity
\eqref{eq:direct-terminal-radius-jets}, with \(Q=\delta\ell^2\), gives
\[
 \delta\log(h_\sigma(x)/Q)
 =-\mathcal A_\sigma(x)+\delta M_\sigma(x),\qquad
 \mathcal A_\sigma(x)=\int_0^x\frac{u\,du}{K_\sigma(u)},\qquad
 \|M_\sigma\|_{C^4}\le C .
\]
The fixed terminal height here is \(z_e/Q=d_f\).

The leading action relation must first give two passages with exactly
the same central height. The assumed margins and compact negative
inverse provide this matching: on the relevant terminal interval,
\((\mathcal A_-^0)'(y)\) is bounded below. At the endpoints of a slightly
enlarged negative radius
interval, the trial passages have central heights bracketing the
positive passage's central height. The order of central heights
(\cref{lem:central-height-order}) gives a matching negative passage.
Implicit differentiation then gives
\begin{equation}
\label{eq:direct-quadratic-action-relation}
 x_-(x)=f(x)+O_{C^3}(\epsilon_3+\delta).
\end{equation}

\emph{Preparing the endpoints.}
Apply \cref{lem:direct-small-height-continuation} with the scales
\eqref{eq:direct-quadratic-prefix-scales},
\(\mathcal U_0=x_\sigma^2/2\) and the constant entrance height factor
\(\zeta_0=d_f\). Both radius intervals are compact and positive, and
the assumed drift margins persist for small \(\epsilon_3\).
Choose the common-cone height \(Z_1=Z_0\ge1\). The endpoints are
\[
 I_\sigma(x_\sigma)=b_\sigma-\tfrac12EA_\sigma x_\sigma^2
                        +Ee_\sigma(x_\sigma),\qquad
 A_\sigma\asymp1,\quad
 \|e_\sigma\|_{C^3}\le C\delta(1+|\log\delta|),
\]
with \(A_-/A_+\le\exp(2C_{\rm cone}\sqrt{z_*}+o(1))\).
Here \(E=\ell^2(z_*/(\ell^2Z_0))^{-2q_2}\); by
\eqref{eq:direct-tail-amplitude} with \(t=\ell^2Z_0\),
\(E\le z_*^{1/4}\ell^{3/2}\), since \(Z_0\ge1\) and
\(|q_2|<1/8\) for small \(\delta\).

Whenever the interval contains a central height in \(\cH_p\), the
endpoints are within \(O(E)\) of admissible endpoints and therefore
lie in \(\widetilde W_\pm\), by \cref{lem:exterior-multiplier-new}.
Composing with the exterior transition gives a Taylor remainder of
\(O(E)\) after division by \(E\). Inserting
\eqref{eq:direct-quadratic-action-relation} gives
\begin{equation}
\label{eq:direct-quadratic-completed-model}
 \frac{2\Delta(x)}{EA_+}
 =D+x^2-\lambda f(x)^2+R_\Delta(x),\qquad
 \|R_\Delta\|_{C^3}
       \le C\{\epsilon_3+\delta(1+|\log\delta|)+E\},
\end{equation}
where \(D\) depends only on the parameters and
\(0<\lambda_0\le\lambda\le\bar\lambda<1\).

\emph{Controlling the crossing direction.}
Apply \cref{lem:direct-derivative-ratio-count} on \(r_0\le x\le R\) to
\(\mathcal F(x)=2\Delta(x)/(EA_+)\), with \(\phi=x^2\), \(\psi=f^2\),
and the \(\lambda\in[\lambda_0,\bar\lambda]\) and remainder
\(R_\Delta\) of \eqref{eq:direct-quadratic-completed-model}; put
\(\mathscr Q=\mathcal F'/(2x)\). Here \(\phi'=2x\ge2r_0\),
\(\phi''=2\), and \(\psi'/\phi'=R_1\) by
\eqref{eq:direct-quadratic-critical-identity}.

At \(b=0\), equal actions give \(y=x\) and \(R_1=1\).
Uniform smooth dependence on the compact drift domains gives
\(|R_1-1|\le Cb\). Fix \(b_*>0\) so small that
\(1-\lambda R_1\ge(1-\bar\lambda)/2\) for \(0\le b\le2b_*\).
The second case of the lemma gives \(\mathscr Q>0\), and hence
\(\Delta_x>0\) in this range. For \(b\ge b_*\),
\cref{lem:direct-equal-action-comparison} gives \(R_1'\ge c_*>0\),
so the first case of the derivative-ratio lemma gives
\begin{equation}
\label{eq:direct-quadratic-weighted-derivative}
 \mathscr Q'<0 .
\end{equation}
These bounds and cuts are fixed before \(\delta,\epsilon_3\) are reduced
to meet the remainder threshold.

To express these signs in the common logarithmic coordinate, use the
positive central-height formula:
\[
 \frac{h_x}{h}=-\frac{\mathcal A_+'(x)}\delta+M_+'(x)<0.
\]
Thus \(x_T>0\) for every such coordinate \(T\). The contraction
interval of large central heights and small radii has
\(\mathcal L<0\), including the cut \(x=r_0\). For \(b\le2b_*\),
\(\Delta_x>0\) and \(h_x<0\) give \(\mathcal L<0\) on the
remaining interval as well. For \(b\ge b_*\), a zero of
\(\mathcal L\) in the compact-radius part has \(\mathcal F'=0\).
By \eqref{eq:direct-quadratic-weighted-derivative},
\[
 \Delta_{xx}=EA_+x\mathscr Q'<0,\qquad
 \Delta_{TT}=\Delta_{xx}x_T^2<0.
\]
With \(m_+=hI_{+,h}>0\), the exact identity
\[
 \Delta_T=-\frac{m_+}{c_T}(\e^{\mathcal L}-1)
\]
therefore gives
\begin{equation}
\label{eq:direct-quadratic-prefix-zero-orientation}
 \mathcal L=0\quad\Longrightarrow\quad
 \mathcal L_T=-\frac{c_T}{m_+}\Delta_{TT}>0.
\end{equation}

Finally, the central heights considered here form one interval in
\(\cH_p\), by \cref{lem:central-height-order}. In increasing \(T\),
the contraction interval and the compact-radius part cover it.
The two estimates apply to the same displacement and include their
shared cut. The compact-radius estimate also includes the terminal cut
\(x=R\), where the endpoints lie in \(\widetilde W_\pm\) and the
exterior map is smooth.
Thus every zero of \(\mathcal L\) on this interval satisfies
\(\mathcal L_T>0\). The second case of
\cref{cor:ordered-cover-zero-count} gives the two-zero bound.
\end{proof}

The compact comparison ends before the positive drift vanishes.
The next lemma gives a strict overlap with the boundary-flow estimate
near that root. It then combines both estimates for the same
logarithmic slope ratio.

\begin{lemma}[Joining the early and root intervals]
\label{lem:prefix-root-cover}
Under \eqref{eq:direct-quadratic-prefix-scales}, suppose \(R\) ranges
over a compact positive interval
and \(K_+^0\) has a fixed positive margin on \(0\le x\le R\).
The negative equal-action radii then lie in a fixed compact interval
where \(K_-^0\) has a positive margin, as required in
\cref{lem:direct-quadratic-prefix}.

Suppose further that \(K_+\) has a first positive root \(\alpha\) in a
fixed compact positive interval. Choose \(0<w_a<R<\alpha\), with fixed
positive lower bounds for \(w_a\) and \(R-w_a\), and a fixed small
\(Y_f>0\) such that the inverse boundary-flow coordinate is defined on
\([w_a,\alpha]\) up to \(Y_f\). Fix a small \(d_f>0\), put
\(z_e=\delta\ell^2d_f\), and use a logarithmic height coordinate \(T\)
as in \eqref{eq:direct-source-clock-convention}. For sufficiently small
\(\delta,\epsilon_3\), the early heights in \(\cH_p\)
\[
 h\ge z_e,\qquad\hbox{or}\qquad
 h<z_e,\quad |\eta_+(z_e;h)|/\ell\le R
\]
form an initial interval in \(T\); every zero of \(\mathcal L\) there
satisfies \(\mathcal L_T>0\). Every remaining height has its positive
projected state at \(z=\ell^2Y_f\) in \([w_a,\alpha)\), with a strict overlap at the
early cut. If every critical point of \(\mathcal L\) in this root range
satisfies \(\mathcal L_{TT}<0\), then \(\Delta_p\) has at most three distinct
zeros on \(\cH_p\); if \(\mathcal L_T>0\) at every zero there, it has at
most two.
\end{lemma}

\begin{proof}
The fixed radius and drift margins bound the positive terminal actions.
To obtain a compact negative inverse, note that
for \(\bar c\ge0\), the action of \(a+by+\bar c y^2\)
increases to infinity; for \(\bar c<0\), it diverges at the first
positive root. Its inverse at each bounded positive action therefore
lies strictly inside the positive-drift interval. Smooth dependence
at a larger fixed action and compactness of the coefficient set give
the claimed compact interval. Thus
\cref{lem:direct-quadratic-prefix} applies to the early interval.

For every remaining height, the positive graph hits \(R\) below
\(Y=\delta d_f\). By \eqref{eq:direct-projection-energy}, its projected
state there lies between \(R-C\delta d_f>w_a\) and \(R<\alpha\).
\Cref{rem:direct-projected-continuation} keeps it increasing and below
\(\alpha\), and continues the inverse projection to \(Y_f\), also at a double
root. The same strict bounds hold at the early cut and give an overlap.
\Cref{lem:central-height-order} orders the two intervals, and
\cref{cor:ordered-cover-zero-count} gives the stated zero bounds.
\end{proof}
